\section{Studi Kasus: Membangun Halaman Profil dengan HTML5}

Studi kasus ini mengintegrasikan seluruh konsep HTML yang telah dipelajari dalam bab ini untuk membangun halaman profil personal yang lengkap dan semantik. Halaman profil adalah jenis halaman web umum yang menampilkan informasi tentang individu, termasuk nama, foto, biografi singkat, daftar keahlian, dan link ke media sosial. Melalui studi kasus ini, Anda akan melihat bagaimana elemen-elemen HTML seperti struktur dokumen, heading, paragraf, daftar, gambar, link, dan elemen semantik bekerja bersama untuk menciptakan halaman web yang bermakna, aksesibel, dan terstruktur dengan baik.

Perencanaan struktur halaman profil dimulai dengan mengidentifikasi bagian-bagian utama yang perlu ditampilkan. Bagian \code{<header>} akan berisi nama lengkap sebagai \code{<h1>} dan tagline singkat. Navigasi dalam \code{<nav>} menyediakan link internal ke bagian-bagian halaman. Bagian \code{<main>} berisi konten utama: foto profil dalam \code{<figure>} dengan \code{<figcaption>}, biografi dalam beberapa \code{<p>}, daftar keahlian dalam \code{<ul>}, dan riwayat pendidikan dalam \code{<dl>}. Bagian \code{<footer>} menutup halaman dengan informasi kontak dan copyright. Pendekatan ini mengikuti pola HTML5 semantik yang meningkatkan aksesibilitas dan SEO.

Implementasi memperhatikan praktik terbaik yang telah dibahas. Atribut \code{alt} pada gambar memberikan deskripsi yang berguna bagi pengguna pembaca layar. Link eksternal ke media sosial menggunakan \code{target="\_blank"} dengan \code{rel="noopener noreferrer"} untuk keamanan. Hierarki heading \code{h1} ke \code{h2} ke \code{h3} mengikuti struktur logis tanpa melompati level. Meta viewport dalam \code{<head>} memastikan halaman responsif di perangkat mobile. Charset UTF-8 menangani karakter internasional dengan benar.

\begin{contoh}
\textbf{Langkah-langkah Membangun Halaman Profil:}

\begin{enumerate}[leftmargin=*]
  \item \textbf{Setup Dokumen:} Buat struktur HTML5 dengan \code{doctype}, \code{html lang="id"}, \code{head} (berisi charset, viewport, title, meta description), dan \code{body}.
  
  \item \textbf{Header:} Tambahkan \code{<header>} dengan \code{<h1>} untuk nama dan \code{<p>} untuk tagline profesional.
  
  \item \textbf{Navigasi:} Buat \code{<nav>} dengan unordered list berisi link anchor ke bagian \code{\#tentang}, \code{\#keahlian}, dan \code{\#kontak}.
  
  \item \textbf{Konten Utama:} Gunakan \code{<main>} sebagai container. Di dalamnya:
  \begin{itemize}
    \item \code{<section id="tentang">} dengan \code{<h2>} dan \code{<figure>} berisi foto dengan \code{<figcaption>}, diikuti paragraf biografi.
    \item \code{<section id="keahlian">} dengan \code{<h2>} dan \code{<ul>} daftar keahlian teknis.
    \item \code{<section id="pendidikan">} dengan \code{<h2>} dan \code{<dl>} untuk riwayat pendidikan.
  \end{itemize}
  
  \item \textbf{Sidebar (Opsional):} Tambahkan \code{<aside>} dengan link media sosial menggunakan ikon atau teks dengan atribut keamanan yang benar.
  
  \item \textbf{Footer:} Tutup dengan \code{<footer>} berisi email (\code{mailto:}), telepon (\code{tel:}), dan copyright.
\end{enumerate}
\end{contoh}

\begin{webcode}[language=HTML, caption={Kode Lengkap Halaman Profil HTML5}]
<!DOCTYPE html>
<html lang="id">
<head>
  <meta charset="UTF-8" />
  <meta name="viewport" 
        content="width=device-width, initial-scale=1.0" />
  <meta name="description" 
        content="Profil profesional Andi Wijaya - 
                 Web Developer dan Desainer UI/UX" />
  <title>Andi Wijaya - Web Developer</title>
</head>
<body>
  <header>
    <h1>Andi Wijaya</h1>
    <p>Web Developer & Desainer UI/UX</p>
    <nav>
      <ul>
        <li><a href="#tentang">Tentang</a></li>
        <li><a href="#keahlian">Keahlian</a></li>
        <li><a href="#kontak">Kontak</a></li>
      </ul>
    </nav>
  </header>
  
  <main>
    <section id="tentang">
      <h2>Tentang Saya</h2>
      <figure>
        <img src="images/profil.jpg" 
             alt="Foto profil Andi Wijaya" 
             width="200" height="200" />
        <figcaption>Andi Wijaya - Web Developer</figcaption>
      </figure>
      <p>Saya adalah web developer dengan pengalaman 5 tahun 
         dalam membangun aplikasi web modern. Spesialisasi 
         saya meliputi <strong>front-end development</strong> 
         dengan HTML, CSS, dan JavaScript.</p>
      <p>Saya percaya bahwa <em>user experience</em> adalah 
         inti dari setiap produk digital yang sukses.</p>
    </section>
    
    <section id="keahlian">
      <h2>Keahlian Teknis</h2>
      <ul>
        <li>HTML5 & CSS3</li>
        <li>JavaScript (ES6+)</li>
        <li>React & Vue.js</li>
        <li>Responsive Web Design</li>
        <li>Git & Version Control</li>
      </ul>
    </section>
    
    <section id="pendidikan">
      <h2>Riwayat Pendidikan</h2>
      <dl>
        <dt>S1 Teknik Informatika</dt>
        <dd>Universitas Indonesia, 2018-2022</dd>
        <dt>Bootcamp Web Development</dt>
        <dd>Codecademy Pro, 2023</dd>
      </dl>
    </section>
  </main>
  
  <aside>
    <h3>Media Sosial</h3>
    <ul>
      <li>
        <a href="https://linkedin.com/in/andiwijaya" 
           target="_blank" 
           rel="noopener noreferrer">LinkedIn</a>
      </li>
      <li>
        <a href="https://github.com/andiwijaya" 
           target="_blank" 
           rel="noopener noreferrer">GitHub</a>
      </li>
    </ul>
  </aside>
  
  <footer id="kontak">
    <h2>Kontak</h2>
    <p>Email: <a href="mailto:andi@example.com">andi@example.com</a></p>
    <p>Telp: <a href="tel:+628123456789">0812-3456-789</a></p>
    <p><small>&copy; 2026 Andi Wijaya. Hak Cipta Dilindungi.</small></p>
  </footer>
</body>
</html>
\end{webcode}

\begin{catatan}
\textbf{Poin Kunci dari Studi Kasus:}
\begin{itemize}[leftmargin=*]
  \item Struktur semantik HTML5 meningkatkan aksesibilitas dan SEO.
  \item Atribut \code{alt} pada gambar wajib ada untuk screen reader.
  \item Link eksternal selalu gunakan \code{rel="noopener noreferrer"}.
  \item Anchor links (\code{\#id}) memudahkan navigasi dalam halaman.
  \item Meta viewport dan charset adalah elemen wajib untuk web modern.
\end{itemize}
\end{catatan}

